\begin{pastexam}[label={px:gk3-114-ds-02}]{高考三級資訊處理 114 年 第二題}
（申論；20 分）

有一個三維整數陣列 A[3][6][8]，每個元素占用 4 個記憶體空間，每個記憶體空間均有位址。該陣列在儲存至
記憶體時，會先被轉換為一維陣列的形式儲存。下列位址皆為十進位，已知 A[0][1][2]的記憶體位址為 2040，
A[1][4][5]的位址為 2340。請問陣列 A 在記憶體中的儲存方式為何？是以列為主（row-major）還是以行為主
（column-major）？（10 分）\\
請計算 A[1][5][3]在記憶體中的位址為何？（10 分）
\end{pastexam}

\begin{pastsoln}
答案：以列為主（row-major）；A[1][5][3] 的位址是 \textbf{2364}。

原卷沒有印小題編號，兩問各 10 分。 宣告 \texttt{A[3][6][8]} 是 C 的寫法，三個下標的範圍是 $0..2$、$0..5$、
$0..7$。 設 $A[0][0][0]$ 的位址是 $\alpha$，每個元素 $4$ 個單位（資結書第~4.1~節）：
\begin{itemize}
\item 以列為主：最右邊的下標變得最快。 $A[i][j][k]$ 前面有 $i$ 個 $6 \times 8$ 的平面、$j$ 列各 $8$ 個、
再 $k$ 個，位址 $\alpha + 4(48i + 8j + k)$。
\item 以行為主：最左邊的下標變得最快。 $A[i][j][k]$ 前面有 $k$ 個 $3 \times 6$ 的平面、$j$ 行各 $3$ 個、
再 $i$ 個，位址 $\alpha + 4(18k + 3j + i)$。
\end{itemize}

\textbf{判斷儲存方式}。 $\alpha$ 不知道，所以比兩個元素的位址差：實際的差是 $2340 - 2040 = 300$ 個單位，
也就是 $300 / 4 = 75$ 個元素。
\begin{itemize}
\item 以列為主：$A[0][1][2]$ 的偏移是 $48 \cdot 0 + 8 \cdot 1 + 2 = 10$，$A[1][4][5]$ 是
$48 \cdot 1 + 8 \cdot 4 + 5 = 85$，差 $85 - 10 = 75$ 個元素，和實際相符。
\item 以行為主：$A[0][1][2]$ 的偏移是 $18 \cdot 2 + 3 \cdot 1 + 0 = 39$，$A[1][4][5]$ 是
$18 \cdot 5 + 3 \cdot 4 + 1 = 103$，差 $64$ 個元素，也就是 $256$ 個單位，不是 $300$。
\end{itemize}
只有以列為主符合，所以是 \textbf{以列為主（row-major）}。 這時起始位址
$\alpha = 2040 - 4 \cdot 10 = 2000$。

\textbf{A[1][5][3] 的位址}。 偏移 $48 \cdot 1 + 8 \cdot 5 + 3 = 91$ 個元素，
\[
\mathrm{loc}(A[1][5][3]) = 2000 + 4 \cdot 91 = 2000 + 364 = \mathbf{2364}.
\]
檢查：從 $A[1][4][5]$ 往後數，$A[1][4][6]$、$A[1][4][7]$、$A[1][5][0]$、$\ldots$、$A[1][5][3]$ 共 $6$ 個
元素，$2340 + 4 \cdot 6 = 2364$，一致。
\end{pastsoln}
